\documentclass{article}
\usepackage{amsmath,amssymb,hyperref}
\title{Generated native state towers from the original bicategorical factory}
\date{15 November 2026}
\begin{document}\maketitle
\paragraph{Status.} New construction target for independent review. This is a possible global iteration route on the tower generated from the original bicategorical process $K_0$. It does not assert a state endofunctor on all process tricategories or bypass the witness equations of the higher dialgebra construction.

\paragraph{Independently constructed normalized levels.} For each finite $n$, let
\[
 T_n=I_{S^n}(K_0;m_n),\qquad
 m_n(X_1,\ldots,X_n)=X_n\otimes\cdots\otimes X_1.
\]
These global native state tricategories are already constructed directly from $K_0$ using the reviewed state-base-change theorem. Their objects are the original interfaces. A machine retains $n$ state coordinates and a $K_0$ body. A native two-cell retains $n$ state maps and one $K_0$ process square. A native three-cell retains $n$ state two-cells satisfying the original process-square equality. The unit-state embedding $\eta_n:K_0\to T_n$ inserts $n$ unit states with all actual unitor comparisons. The seed $T_0$ uses the terminal state bicategory $S^0$ and is equivalent to $K_0$ through the unit-action comparisons. This does not assert that $\eta_n$ is full on all native maps. No level is defined by first assuming the next level exists.

\paragraph{Raw next-level native hom.} Suppose a previous raw level has been identified homwise with $T_n$, preserving its unit-state pure embedding. For fixed interfaces $A,B$, define the outer-state trihomomorphisms
\[
 P_A:S\longrightarrow T_n,\quad X\longmapsto JX\otimes A,
 \qquad Q_B:S\longrightarrow T_n,\quad X\longmapsto JX\otimes B,
\]
by composing the original pure-interface action with $\eta_n$. On state arrows these are actual $n$-unit-state pure machines. The reviewed higher native dialgebra construction therefore defines its raw next-level hom bicategory, retaining the outer state map and the \emph{actual} inner process-three-cell witness. This requires these individual trihomomorphisms, not an unproved complete higher bimodule-action package on arbitrary tricategories.

\paragraph{Inductive normalization target.} Its objects have outer state $X_1$ and an inner $n$-state machine with coordinates $X_2,\ldots,X_{n+1}$. Normalize the body by the specified action equivalence to a $K_0$ body on $X_{n+1}\otimes\cdots\otimes X_1$. An outer native arrow has state map $h_1$ and an inner native arrow between composites with the $n$-unit-state boundary machines. At each inner coordinate use the chosen unitors to normalize the unit-framed map to $h_i$. The remaining process square has boundary $m_{n+1}(h_1,\ldots,h_{n+1})$.

An outer native three-cell has $\sigma_1$ and an inner witness. The inductive presentation of $T_n$ resolves that witness into $\sigma_2,\ldots,\sigma_{n+1}$ and the actual $K_0$ equation. After unit-frame transport this is exactly the native equation for $m_{n+1}(\boldsymbol\sigma)$. No coordinate is forgotten, and equality of images under tensor is never used to infer equality of tuples. This proposes a homwise biequivalence to $T_{n+1}(A,B)$. The inverse inserts the same unit frames; its units/counits are the specified unitor adjunctions. Composition comparison is induced by their pseudonaturality. More directly, this is the heterogeneous two-level normalization with inner state bicategory $S^n$ acting through $m_n$ and outer state bicategory $S$. An inner native witness is by definition a tuple of $n$ state two-cells and one $K_0$ equation; no hypothetical arbitrary higher witness is recursively simplified. The reviewed two-level proof therefore applies to $(S,S^n)$ with its specified action comparison, including boundary-state and process-witness naturality.

\paragraph{Global assembly by explicit normalization.} Rather than postulate a complete action on arbitrary process tricategories, construct global assembly on these raw next-level homs by transporting the already established $T_{n+1}$ structure along the specified homwise biequivalences. Choose coherent biadjoint biequivalence data for each hom normalization, not merely arbitrary quasi-inverse object maps. The global transport uses those inverse pseudofunctors, pseudonatural equivalences and their coherence modifications. This is stronger data than citing an object correspondence; its existence and use must be checked as part of the transport audit. The transported composition is: normalize both arrows, compose in $T_{n+1}$, and insert unit frames. Its associators, unitors and modifications are the corresponding pasted cells of $T_{n+1}$ and those equivalences. Their axioms follow after applying the hom equivalences, which are faithful on the relevant hom two-cells; this is \emph{not} the nonfaithful projection forgetting process witnesses.

The result must also be compared with the intended raw chronological operation, not merely equipped with an arbitrary transported product. That operation retains the outer state pair and every inner coordinate pair, runs the first original body while holding the other coordinates, then the second. On witnesses it pastes the actual inner witnesses and the pure boundary comparisons. Normalization yields exactly the coordinatewise state pairs and chronological body of $T_{n+1}$; the comparison is the same matrix-transposition/action equivalence as the reviewed global finite joining theorem. The comparison on every native witness is given explicitly in the following calculation. It does not require an inverse effectful interchange.

\paragraph{The chronological comparison on actual witnesses.} Write a raw machine as $(X,\mathbf A,f)$ and its normalization as the body $\rho_B f\rho_A^*$, where $\rho$ is the specified action equivalence from $L_{m_n\mathbf A}L_X$ to $L_{m_{n+1}(X,\mathbf A)}$. For a second machine $(Y,\mathbf B,g)$, the raw chronological composite retains
\[
 (X\otimes Y,(A_i\otimes B_i)_i).
\]
Normalization of its held-first/held-second body and composition of the normalized bodies differ by the specified strong monoidal comparison of $m_n$ and the original action's middle-coordinate permutation. Define $\kappa$ by precisely that structural paste, including chosen adjoint inverses in body transport, not by an arbitrary equivalence between endpoints.

For native squares $\alpha,\beta$, expand both sides of $\kappa$ naturality. Strip their unit-state frames, apply $\alpha$ first and $\beta$ second, and close the same coordinatewise state boundary. Pseudonaturality of the action comparisons at these actual squares identifies the pastes. For native witnesses $\Gamma,\Delta$, the normalized state components are the same coordinatewise two-cells $\sigma_i\otimes\tau_i$ with their monoidal compositors. Substituting the two actual $K_0$ native-square equations into the chronological paste proves the required witness equation. No faithfulness of $m_n$ and no inverse effectful $\chi$ is used.

For three machines, the associativity diagram for $\kappa$ is the monoidal-composition coherence of $m_n$ and the original action associativity/middle-permutation coherence, evaluated on the fixed chronological word $f;g;h$. The unit diagrams use the unit comparison of $m_n$ and the action unit modifications. Each needed native cell first has its $K_0$ admissibility established by these equations. Equations between admitted modifications can then be checked after normalization in $T_{n+1}$, whose tuple-state projection is faithful on native three-cells. This projection retains the entire tuple; it is not the tensor map and not the process-witness-forgetting map of an arbitrary process tricategory. Thus the specified chronological assembly agrees coherently with the transported global structure.

\paragraph{Typed homwise transport audit.}
Here is the many-interface transport used above. Write $\mathcal R(A,B)$ for the raw hom bicategory and $\mathcal T(A,B)=T_{n+1}(A,B)$. Fix biadjoint biequivalences
\[
 N_{AB}:\mathcal R(A,B)\rightleftarrows\mathcal T(A,B):R_{AB},
 \quad u_{AB}:1\simeq R_{AB}N_{AB},\quad
 e_{AB}:N_{AB}R_{AB}\simeq1,
\]
including their coherent triangle modifications. For chronological composition
$m_{ABC}:\mathcal T(A,B)\times\mathcal T(B,C)\to\mathcal T(A,C)$ set
\[
 \widetilde m_{ABC}=R_{AC}m_{ABC}(N_{AB}\times N_{BC}),
 \qquad\widetilde I_A=R_{AA}I_A.
\]
These are pseudofunctors and specified objects, respectively. The comparison
$N_{AC}\widetilde m_{ABC}\simeq m_{ABC}(N_{AB}\times N_{BC})$
is the counit $e_{AC}$ evaluated at that composite; the identity comparison is $e_{AA}(I_A)$. Thus its direction and its source hom bicategory are fixed.

To construct the associator, first conjugate the associator of $\mathcal T$ by these counit comparisons: on the left nested composite apply $e$ at the outer composite and then $m(e,1)$ at the inner composite, apply the $\mathcal T$ associator, and reverse the corresponding two structural comparisons on the right. This is a pseudonatural equivalence between the $N$-images of the two raw triple-composition pseudofunctors. Lift it using $R$ and the unit $u$ at these two raw functors. Equivalently this uses the biequivalence induced by postcomposition with $N$ on the appropriate pseudofunctor bicategory. Unit equivalences are obtained in the same way. No native process witness is inverted in these constructions; the inverses are those of the specified equivalence data.

The pentagonator and the three tricategorical unit modifications are obtained by the corresponding conjugated modifications of $\mathcal T$, with the comparison modifications for $N,R,u,e$ inserted. In particular they are not declared identities merely because their objectwise endpoints have the same normalized state. Every coherence equation is checked after postcomposition with $N$ and these comparison identifications. The residual pastings are the actual tricategory equations of $\mathcal T$, pseudonaturality, and the coherent triangle modifications of the biadjoint biequivalences. Postcomposition is locally fully faithful on modifications between the fixed pseudonatural transformations, so equality reflects. This argument concerns the relevant hom categories of modifications; it does not assert that the object function of $N$, or the tensor of the state coordinates, is injective.

The resulting homwise normalization, with its counit composition/identity comparisons and the just constructed coherence modifications, is a trihomomorphism and a local biequivalence. This establishes the global transport used here, not an arbitrary transfer of additional higher action or symmetric monoidal structures. Those must still be compared with their specified operations, as done above for chronological assembly.

\paragraph{Transport background and attribution.}
Gurski's \emph{Biequivalences in tricategories},
\href{https://arxiv.org/pdf/1102.0979v1}{arXiv:1102.0979v1},
Theorem 3.2 supplies biadjoint biequivalence data for biequivalences of bicategories; Theorem 4.5 gives the general tricategorical result. Theorem 5.1 treats transport of a monoidal bicategory along a biequivalence. The hom-indexed construction above is the explicit many-interface transport needed here, not a claim that the one-object monoidal theorem alone proves our native normalization or chronological comparison. No novelty claim is made for this transport mechanism.

\paragraph{Insertion and joining.} At the normalized levels, insertion of a unit coordinate and joining consecutive coordinates are the strong symmetric monoidal state-base maps already studied. Their global comparisons retain unit-state whiskers and conjugate process bodies. Transporting these maps across the raw-level biequivalences gives candidate unit insertion and finite joining on the generated raw tower. The associativity/unit comparisons are the reviewed block-joining comparisons, now with normalization equivalences included. This would give coherent iteration and joining on the generated finite tower, not by itself a pseudomonad on an unspecified four-dimensional category of all process tricategories.

\paragraph{Honest boundary.} The proof target is all finite levels generated from the actual original relative factory, with arbitrary admitted native maps and process witnesses at those levels. It is not an operational special-case family. Its limitation is ambient: arbitrary process tricategories not generated in this way still need their own higher action/assembly hypotheses and proof. A general evaluation universal property, infinite-depth completion and arbitrary shared-diagonal synchronization remain separate. The point is to derive the next generated level from independently constructed normal forms while preserving all raw native cells, rather than demand a stronger ambient theorem before making actual iteration available.
\end{document}
